\chapter{几何力学}\label{chap:geometric-mechanics}
\begingroup
\renewcommand{\thefootnote}{}
\footnotetext{\textcopyright\ Springer International Publishing Switzerland 2014.
  L. Godinho and J. Natário, \emph{An Introduction to Riemannian Geometry},
  Universitext. DOI \href{https://doi.org/10.1007/978-3-319-08666-8_5}{10.1007/978-3-319-08666-8\_5}.}
\endgroup

力学是研究运动的科学, 基本上始于 Galileo 及其革命性的实证方法. 第一个精确的数学表述由 Newton 在 \emph{Philosophiae Naturalis Principia Mathematica} 中给出. 该书于 1687 年首次出版, 内容丰富, 其中包括对行星绕太阳运行的椭圆轨道的解释. Newton 的思想由多位数学家发展并推广, 包括 Euler, Lagrange, Laplace, Jacobi, Poisson 和 Hamilton. 特别是, 天体力学达到了极高的精度: 例如, 1846 年发现 Neptune, 正是因为需要解释 Uranus 的观测轨道与理论预测之间的偏差.

本章利用 Riemann 几何给出 Newton 力学的几何表述. 如第~\RGref{sec:mechanical-systems}~节所解释, 这之所以可行, 是因为任何力学系统的动能都在其位形空间上给出一个 Riemann 度量. 位形空间是一个可微流形, 其点表示系统可能的位形. 第~\RGref{sec:holonomic-constraints}~节说明\textbf{完整约束}如何产生非平凡的力学系统: 这类约束迫使系统沿位形空间的子流形运动. 第~\RGref{sec:rigid-body}~节详细研究一个特别重要的例子, 即\textbf{刚体}. 第~\RGref{sec:nonholonomic-constraints}~节讨论限制速度而不是位形的\textbf{非完整约束}.

第~\RGref{sec:lagrangian-mechanics}~节给出力学的 Lagrange 表述, 其中运动轨迹是使作用量积分取极值的曲线. 本节也讨论 \textbf{Noether 定理}, 它把守恒律与对称性联系起来. 第~\RGref{sec:hamiltonian-mechanics}~节说明与之对偶的 Hamilton 力学表述, 其中运动轨迹由余切丛中的特定流得到. 第~\RGref{sec:completely-integrable-systems}~节用它建立\textbf{完全可积系统}的理论, 这类系统的动力学特别简单. 第~\RGref{sec:symmetry-and-reduction}~节将 Hamilton 形式体系推广到辛流形和 Poisson 流形, 并讨论这些流形在适当对称性下的约化.
